\chapter{Introducción}
\hyphenation{In-te-re-se}

Las nubes de puntos 3D son una representación de una realidad, como un objeto o un terreno, mediante un conjunto de puntos en el espacio tridimensional. Cada uno de esos puntos está caracterizado por sus coordenadas espaciales, así como otras propiedades como el color o la intensidad. Entre las aplicaciones más importantes de la elaboración de estas nubes están la cartografía, la elaboración de modelos tridimensionales en CAD de piezas fabricadas, la metrología, la arqueología o la conducción autónoma.
La tecnología predominante para la captura de estos puntos es la LiDAR (Light Detection And Ranging).

Una de las tareas más comunes al trabajar sobre nubes de puntos 3D es la búsqueda de vecindades alrededor de un punto dado. Una vecindad se define como el conjunto de puntos que están dentro de una figura geométrica determinada en el espacio con radio $R$ y con el centro en el punto para el cual se están buscando vecinos. Esta combinación de figura geométrica y de radio se conoce como kernel de búsqueda. Otro método alternativo de análisis de vecindades es buscar los $K$ puntos más cercanos al elemento de partida, técnica llamada KNN (\textit{k nearest neighbors}). Cabe destacar, eso sí, que el presente trabajo se centra exclusivamente en el primer enfoque (búsquedas de geometría de kernel), quedando las búsquedas KNN fuera de alcance.

Un gran problema a nivel de hardware de este proceso sobre nubes de puntos es que estos datos se recogen sin ningún tipo de estructura ni orden basado en localidad. Este fenómeno se junta con el hecho de que estas nubes pueden tener millones de puntos almacenados. Por eso, se necesitan métodos eficientes para acelerar este proceso, intrínsecamente costoso, de acceso a puntos cercanos.

Los \textit{Octrees} y los \textit{KDTrees} son estructuras jerárquicas ampliamente utilizadas para lidiar con este problema de ineficiencia. Consisten en dividir la nube de puntos de manera recursiva en zonas próximas en el espacio físico. Los \textit{Octrees} son estructuras donde cada nodo se divide exactamente en 8 hijos, dividiendo cada dimensión a la mitad (solo se usa en espacios tridimensionales). Los puntos pertenecen al nodo hoja que los contiene en su interior. Por otra parte, los \textit{KDTrees} son árboles binarios que dividen cada nodo interno solamente según una dimensión. En cada nivel se elige un eje y se traza un hiperplano que divide a los puntos en dos nodos. De esta forma, este árbol es generalizable para cualquier espacio de dimensiones.

El rendimiento de estas estructuras depende directamente de la distribución de los puntos. Por eso es muy útil la aplicación de técnicas de reordenamiento espacial basadas en curvas de llenado de espacio (SFC). Su objetivo es mejorar la coherencia entre la localidad espacial de los puntos y su disposición en memoria para optimizar los accesos en el cálculo de puntos vecinos.

En este trabajo se usará como base el modelo de \textit{Octrees} propuesta por Viñambres et al.\cite{Vinambres},  donde se aplican las reordenaciones espaciales de Morton y Hilbert sobre diferentes \textit{Octrees} implementados en el lenguaje \texttt{C++}. Se construyó tanto una estructura clásica basada en punteros enlazados como una versión basada en almacenamiento en un vector lineal. Sobre estas estructuras, y los algoritmos de búsqueda de vecindades asociados a ellas, se tratará de optimizar todavía más el acceso a puntos cercanos. Para ello se empleará la técnica algorítmica de la poda. 

Esta estrategia, de forma general, consiste en descartar o \textit{podar} regiones del espacio (nodos) o conjuntos de datos para los que matemáticamente se puede demostrar que no satisfacen las condiciones necesarias para ser considerados solución. Su aplicación en este contexto es ser capaz de conocer \textit{a priori} qué partes del \textit{Octree} no pueden contener ningún vecino para una búsqueda dada. Esta técnica es, de hecho, el motivo por el que los \textit{Octrees} (y los KDTrees) son estructuras ideales para almacenar nubes de puntos 3D. Es muy sencillo podar nodos comprobando si sus límites están fuera del alcance del kernel de vecindad. Sin embargo, en este trabajo se explorará la poda a nivel de puntos dentro de un nodo, es decir, se busca ser capaz de descartar puntos pertenecientes a un nodo hoja sin tener que evaluarlos individualmente. De esta forma, se acotaría un rango específico dentro del dominio espacial del nodo y se rechazarían sin necesidad de recorrerlos secuencialmente todos los puntos de la hoja que no perteneciesen a ese rango. Para conseguir esta tarea se realizarán reordenamientos de puntos a nivel de hoja siguiendo diversos criterios basados en coordenadas cartesianas con ejes X, Y, Z como definidores de los puntos, pero también se manejarán sistemas alternativos de representaciones tridimensionales: sistemas esféricos y cilíndricos. Estos sistemas son expansiones a espacios tridimensionales de las coordenadas polares.

Las coordenadas polares son un sistema de coordenadas bidimensional donde cada punto es caracterizado por una distancia radial y un ángulo con respecto del origen. Pueden ser expandidas a un espacio tridimensional mediante dos sistemas ligeramente distintos: las coordenadas esféricas y las cilíndricas.

Los reordenamientos en cuestión consisten en almacenar los puntos dentro de cada nodo terminal ordenados según alguna característica de las mencionadas. Así, se podrán calcular los límites de forma teórica para ese valor que puede tener un punto de forma que satisfaga la restricción de vecindad, dada una búsqueda con un kernel determinado.

Para resumir, esta propuesta de algoritmo y este estudio se pueden descomponer en una serie de objetivos claros, que se definen a continuación:
\begin{enumerate}
    \item  Usar las técnicas de los reordenamientos y las podas para mejorar la eficiencia en las búsquedas de vecinos mediante kernels geométricos sobre nubes de puntos 3D capturadas por sensores LiDAR. Al hablar de eficiencia se hace referencia a la temporal, buscando la mayor rapidez posible. En este trabajo no se tomará como prioridad el ahorro en memoria computacional, aunque como es lógico sí se tendrá en cuenta y se tratará de minimizar su gasto siempre que no genere un impacto negativo en el rendimiento temporal.
    \item Comprobar la validez teórica de la estrategia de los reordenamientos y podas a nivel de hoja en \textit{Octrees}, de forma que se pueda considerar una técnica plausible, más allá de los resultados temporales obtenidos en esta implementación del problema.
    \item Estudiar y analizar el rendimiento de los algoritmos de búsqueda optimizados mediante la citada estrategia frente a los algoritmos de búsqueda de vecindades de partida, de forma que se puedan sacar conclusiones lo más realistas posibles de la efectividad de este proyecto.
    \item Analizar la escalabilidad de la paralelización en los nuevos algoritmos de búsqueda, comparando este factor también con los algoritmos de partida.
\end{enumerate}
